\chapter{Design}
\label{chap:design}
%%%%%%%%%%%%%%%%%%%%

This chapter details the architecture and algorithmic design of the proposed adaptive actuation system.

\section{System Architecture}
The system utilizes a multi-layered software and hardware stack designed to bridge high-frequency physical telemetry with a dynamic optimization engine. At the hardware layer, a modified PRO X2 Superstrike mouse streams continuous 120-level analog depth values at 1\,kHz. This bypasses typical firmware constraints that generally limit reporting to mere 4-bit precision at 65\,Hz and block reading after haptic events. The hardware communicates with the host operating system via an OS-level driver bridge that interfaces high-level prototyping environments with production-grade HID protocols. Sitting above this is a background sensor service that orchestrates and synchronizes the analog telemetry with OS events and other biometrics. Finally, a front-end orchestration client provides the user dashboard, manages recording state, and drives the downstream machine learning pipelines. The overall component topology is shown in Figure~\ref{fig:system_arch}.

\begin{figure*}[htbp]
\centering
\begin{tikzpicture}[
  font=\small\sffamily,
  box/.style={
    draw=black!40, line width=0.7pt,
    rounded corners=5pt, fill=black!5,
    minimum width=2.8cm, minimum height=1.0cm,
    align=center
  },
  arr/.style={
    -{Stealth[length=2.2mm, width=1.4mm]},
    line width=0.9pt, black!65
  },
  biarr/.style={
    {Stealth[length=2.2mm, width=1.4mm]}-{Stealth[length=2.2mm, width=1.4mm]},
    line width=0.9pt, black!65
  },
  lbl/.style={font=\scriptsize\sffamily, fill=white, inner sep=1.5pt, text=black!55},
  bcpp/.style={draw=orange!70,  fill=orange!15,  rounded corners=2pt,
               font=\fontsize{6}{7}\selectfont\bfseries, text=orange!70!black,  inner sep=2pt},
  bpy/.style={ draw=green!55!black, fill=green!12, rounded corners=2pt,
               font=\fontsize{6}{7}\selectfont\bfseries, text=green!50!black,   inner sep=2pt},
  bts/.style={ draw=blue!55,    fill=blue!9,    rounded corners=2pt,
               font=\fontsize{6}{7}\selectfont\bfseries, text=blue!60!black,    inner sep=2pt},
  brust/.style={draw=red!55,    fill=red!9,     rounded corners=2pt,
                font=\fontsize{6}{7}\selectfont\bfseries, text=red!60!black,    inner sep=2pt},
  bnode/.style={draw=violet!55, fill=violet!9,  rounded corners=2pt,
                font=\fontsize{6}{7}\selectfont\bfseries, text=violet!60!black, inner sep=2pt}
]

% ─── Hardware stack (left) ────────────────────────────────────
\node[box] (mouse)  at (0.0, 0.0) {HITS Mouse\\[-1pt]{\scriptsize custom FW}};
\node[box] (driver) at (0.0, 2.1) {HID Driver\\Bridge};
\node[box] (sensor) at (0.0, 4.2) {Background\\Sensor Service};

% ─── Front-end client (centre-top) ───────────────────────────
\node[box, minimum width=3.6cm] (client) at (5.0, 6.1)
      {Front-end\\Orchestration Client};

% ─── Out-of-process backend services ─────────────────────────
\node[box] (decoder) at (3.0, 3.5) {Analog Decoder\\[-1pt]\& Translator};
\node[box] (ml)      at (6.0, 3.5) {ML Manager};
\node[box] (fsm)     at (9.0, 3.5) {FSM Simulator\\[-1pt]\& Optimizer};

% ─── Language badges (top-right corner, all right-aligned) ───
\node[bcpp,  anchor=south east] at ([xshift=2pt,yshift=-4pt]mouse.north east)          {C++};
\node[bcpp,  anchor=south east] at ([xshift=2pt,yshift=-4pt]driver.north east)         {C++};
\node[bcpp,  anchor=south east] at ([xshift=2pt,yshift=-4pt]sensor.north east)         {C++};
\node[bts,   anchor=south east] at ([xshift=2pt,yshift=-4pt]client.north east)         {TS};
\node[bnode, anchor=south east] at ([xshift=-12pt,yshift=-4pt]client.north east)       {Node};
\node[bts,   anchor=south east] at ([xshift=2pt,yshift=-4pt]decoder.north east)        {TS};
\node[bpy,   anchor=south east] at ([xshift=2pt,yshift=-4pt]ml.north east)             {PY};
\node[brust, anchor=south east] at ([xshift=2pt,yshift=-4pt]fsm.north east)            {Rust};

% ─── Arrows ───────────────────────────────────────────────────
% Hardware stack — unidirectional upward
\draw[arr] (mouse)  -- node[lbl, right] {HID}   (driver);
\draw[arr] (driver) -- node[lbl, right] {IPC}   (sensor);

% Sensor service ↔ client (bidirectional diagonal)
\draw[biarr] (sensor.east) --
    node[lbl, above, sloped, pos=0.5] {IPC} (client.west);

% Backend services ↔ client (bidirectional)
\draw[biarr] (decoder.north) to[out=80, in=215]
    node[lbl, left,  pos=0.4] {IPC} (client.south west);
\draw[biarr] (ml.north) --
    node[lbl, right] {IPC} (client.south -| ml.north);
\draw[biarr] (fsm.north) to[out=100, in=330]
    node[lbl, right, pos=0.4] {IPC} (client.east);

\end{tikzpicture}
\caption{System architecture. Telemetry flows upward from a HITS-enabled mouse
with custom firmware through a C++ HID driver bridge and a C++ background sensor
service via IPC to the front-end orchestration client (Node/Electron\,+\,TypeScript).
Three independent out-of-process backend services communicate bidirectionally with
the client via IPC\@: an analog decoder and translator (Python\,+\,TypeScript), a
machine learning manager (Python), and a hardware FSM simulator and optimizer (Rust).}
\label{fig:system_arch}
\end{figure*}

\section{Data Labeling Methodology}
\label{sec:labeling}

Raw analog signals are inherently noisy. Before any label placement, the depth telemetry is smoothed using Modified Akima (MAKIMA) interpolation. MAKIMA was specifically chosen because it preserves signal monotonicity and avoids the severe overshoots characteristic of standard cubic splines, which is critical for accurate boundary detection on a signal with sharp transitions.

\subsection{Algorithmic Bootstrapping}
To alleviate the bottleneck of manual annotation at scale, an automated heuristic pipeline pre-populates boundary candidates before human review. Detection anchors on the digital press/release event pair reported for each click to identify the approximate click window, then walks the smoothed analog signal within that window to locate the true biomechanical boundary:

\begin{itemize}
    \item \textbf{Make Intention:} Starting from the click's reported start, the algorithm walks backwards in time in search of the point of minimal analog depth, treating a run of frames without further improvement as evidence that the true rest position has been passed. This favors the deepest point actually reached over any single noisy sample, since sensor micro-tremors near rest position can otherwise produce spurious local minima that would misplace the boundary.
    \item \textbf{Break Intention:} Detected as the first point, walking forward from the click's reported start, where the (smoothed) rate of depth change falls below a small negative margin---i.e., the finger begins lifting with a speed larger than sensor noise. This typically precedes the reported key release and more accurately reflects the physical start of the finger lift.
\end{itemize}

This heuristic correctly handles the vast majority of clean, well-separated clicks, providing an accurate starting point for human reviewers to correct rather than annotate from scratch. Concrete search parameters and fallback behavior are implementation concerns, described in Section~\ref{sec:labeling_impl}.

\subsection{Iterative Human-in-the-Loop Annotation}
Annotation followed an iterative three-phase strategy designed to maximize label quality while minimizing annotation effort:

\begin{enumerate}
    \item \textbf{Algorithmic bootstrap.} The heuristic pipeline above generated initial boundary candidates for all recordings. Annotators reviewed these within a custom data visualization and labeling interface integrated directly into the front-end orchestration client. The interface renders the raw 1\,kHz analog depth signal alongside its derived velocity, and allows annotators to drag start and end boundary markers independently for the left and right button channels. Visual smoothing was applied on a per-recording basis to assist in distinguishing genuine kinetic events from micro-tremors. Annotators focused strictly on biomechanical intent, preserving deliberate presses including those not registered by the OS, while discarding involuntary micro-tremors and resting-weight drift. Approximately 50\% of the dataset was verified and corrected in this phase.

    \item \textbf{Model-assisted labeling.} Once the manually verified corpus was sufficiently large, a preliminary model was trained and used to auto-label the remainder of the dataset. Segments with a model confidence of 95\% or above were accepted automatically, covering roughly two-thirds of each of the remaining unlabeled recordings.

    \item \textbf{Mixed review.} The remaining clicks (primarily the edge cases) were handled by direct human annotation. A final manual review pass was conducted over all automatically labeled segments to validate quality before inclusion in the training corpus. Non-click segments (after trimming idle background) were manually verified to ensure that no labels were missed, with the help of the make and break probability curves produced by the model to highlight potential false negatives. This final review pass ensured that the dataset was as complete and accurate as possible, with a focus on edge cases and ambiguous signals.
\end{enumerate}

The full extent and breakdown of the resulting annotated corpus is described in Section~\ref{sec:data_acquisition}.

\section{Data Preparation Pipeline}
\label{sec:data_prep}

Raw recordings and their human-verified labels are transformed into normalized arrays shared by all three model architectures. Several design decisions govern this pipeline. First, onset boundary zones are widened by a small tolerance radius to absorb natural annotation imprecision, but are adaptively clamped against adjacent segment boundaries to prevent cross-segment label contamination. Second, derivative channels (velocity, acceleration, jerk, and resting-drift) are derived using zero-phase filtering to avoid introducing causal shift artifacts. Third, idle background periods are trimmed by an activity mask to reduce the dominance of background frames without discarding legitimate negative examples. Fourth, dataset splits are performed at the participant level so no individual's recordings span both training and evaluation. Finally, normalization statistics are fitted on the training partition only, and statistics for kinematic channels are computed over active frames exclusively to prevent background mass from compressing the feature scale. Concrete implementation details---feature extraction windows, normalization constants, and weighting formulas---are described in Section~\ref{sec:ml_backend}.


\section{Click-Intent Sequence Modeling}
\label{sec:modeling}

The click-intent problem is formulated as a per-frame four-class temporal segmentation task over continuous 1\,kHz sensor streams. Each frame is assigned one of four mutually exclusive labels: Background, Make Onset, Click Hold, and Break Onset. Three model architectures are supported, positioned along a speed--accuracy trade-off curve suited to the realities of on-device, per-session deployment.

\subsection{Random Forest}
The Random Forest classifier operates on hand-crafted sliding-window features extracted from a local neighbourhood around each frame. It is the primary production model. Its main strength is training speed: a full session retrains in a few minutes on a consumer CPU. Feature importance scores are directly interpretable, exposing which kinematic cues drive the model's decisions. The main weakness is a bounded temporal horizon limited to a few tens of milliseconds of context, which is sufficient for most click onset detection but may struggle on ambiguous edge cases where the intent signal unfolds over a longer transition. This architecture also fails to capture long-range dependencies and long context of user behavior. The pipeline also requires domain knowledge to design and cannot adapt to latent signal structure that was not anticipated during feature engineering.

\subsection{Gradient Boosted Trees}
The Gradient Boosted Trees model shares the same feature extraction pipeline as the Random Forest but replaces the ensemble learner with sequential boosting, which iteratively focuses on the residuals of prior trees. Its key advantage is support for custom loss functions: a focal loss~\cite{lin_focal_2017} variant down-weights easy background frames, reducing over-confidence on the majority background class. Early stopping against a held-out validation metric limits overfitting. The weaknesses are slower training than the Random Forest, greater hyperparameter sensitivity, and the same fundamental constraint of a short-context hand-crafted feature representation.

A further design option explored for the GBT path is \textbf{asymmetric soft labels}, intended to address two simultaneous sources of asymmetry: natural annotator imprecision at boundary placement, and the directional difference in the cost of timing errors. For the Make Onset, a slightly \emph{early} detection is acceptable---the finger is either already in motion or the signal is flat, so no threshold can trigger a click, and the error has negligible impact on the optimizer's boundary estimate---whereas a \emph{late} detection shifts the perceived onset forward in time, increasing input latency. For the Break Onset the direction of tolerable error is the same: a slightly \emph{early} detection is benign because the signal is still rising or flat, and no break RT threshold can trigger a release on such a portion of the depth trace, whereas a \emph{late} detection holds the click open beyond the true biomechanical release point, directly distorting the break boundary seen by the optimizer. The soft label generator encodes this asymmetry with a split-normal Gaussian centred on each onset frame: a wide left lobe and narrow right lobe for both Make and Break Onsets. This objective is implemented and supported by the training pipeline but was not used for the GBT results reported in this thesis, which are trained on hard labels with the focal loss objective described above.

\subsection{Fully Convolutional Network}
The dilated residual FCN operates directly on raw normalized sensor streams without manual feature engineering. Stacked dilated convolutions grow the receptive field exponentially: with eight dilation stages and a kernel size of 7, the effective receptive field spans approximately 3\,060 frames ($\approx$3\,s at 1\,kHz), allowing the model to condition each frame's prediction on long-range temporal context---pre-click muscle tension ramp-ups and post-click rebound dynamics that fall well outside any practical sliding window. The model is trained end-to-end on a three-channel Gaussian heatmap target, producing soft onset probability curves rather than hard class assignments. The trade-off is higher computational cost: the FCN requires a GPU or significantly longer CPU training time, is less interpretable than tree-based models, and benefits more strongly from larger datasets. It is therefore positioned as the high-accuracy architecture for offline post-processing, while the Random Forest handles latency-sensitive or resource-constrained deployment.

\section{From Classification to Segmentation}
\label{sec:segmentation}

The per-frame classifier produces a four-class probability array over the full recording. Converting these continuous probability outputs into discrete (start, end) click segments requires three sequential design choices: identifying candidate onset regions, electing a single representative frame from each region, and pairing each Make onset region with its corresponding Break onset region.

\paragraph{Candidate run detection.}
Two separate probability channels are extracted from the classifier output: the Make Onset probability and the Break Onset probability. Contiguous regions above a configurable threshold in each channel are identified as candidate runs. Make runs use a hysteresis-based entry/exit threshold pair rather than a single cutoff, so that a single physical press---whose probability hump can dip transiently below threshold due to quantization noise or softmax saturation---is not fragmented into multiple spurious candidates. Break runs are detected with a simpler thresholding scheme, since the Break probability curve is typically sharper and less prone to multi-modal fragmentation. The make and break thresholds are hyperparameters that are jointly optimized with the onset frame election and pairing strategies described below.

\paragraph{Onset frame election.}
Within each candidate run, a single frame must be elected as the representative onset. Three strategies are evaluated, trading off robustness to probability-curve saturation against localization sharpness: \textbf{peak} selects the frame of maximum probability; \textbf{plateau median} instead returns the median frame among all near-peak frames, which is more stable when the curve saturates into a wide plateau; \textbf{steepest rise/fall} selects the inflection point of the probability curve, independent of any plateau. The optimal strategy and its associated threshold values are selected jointly via a Bayesian hyperparameter sweep (Optuna~\cite{optuna_akiba_2019}) over the validation set, optimizing the Intersection over Union (IoU) between predicted and ground-truth click segments.

\paragraph{Make--Break pairing.}
Given a set of elected Make and Break onset frames, a final matching step assigns each Make to exactly one Break. Three pairing strategies of increasing sophistication are considered. A \textbf{greedy first-fit} baseline pairs each Make with the first eligible Break in its search window. A globally optimal \textbf{assignment-based} strategy instead poses pairing as a minimum-cost bipartite matching problem over Make and Break confidence, guaranteeing no Break is claimed by more than one Make at the cost of a less scalable search. A \textbf{trained pairing scorer} further replaces the raw confidence-based cost with a learned classifier over pair-level features (confidences, timing, and local competition), targeting the cases where confidence alone is not a reliable indicator of the correct pairing. Because a globally optimal or learned strategy is not guaranteed to outperform the simpler baseline on every dataset, the three are compared empirically on the validation set, and the best-performing strategy is selected and stored alongside the model weights so inference uses identical pairing logic. Algorithmic and training specifics for each strategy are described in Section~\ref{sec:segmentation_impl}.

\section{Parameter Optimization}
\label{sec:optimizer}

After click-intent boundaries are extracted, the optimizer estimates the hardware threshold pair that best reproduces those boundaries under firmware-faithful trigger simulation.

\subsection{Firmware State Machine Simulation}
Each candidate threshold pair is evaluated with a frame-accurate simulation of firmware trigger logic.

In \textbf{simple-threshold mode}, the simulator is a 2-state machine: press is emitted when the make signal crosses the make threshold upward, and release is emitted when the break signal crosses the break threshold downward.

In \textbf{Rapid Trigger mode}, the simulator is a 3-state machine (Idle $\to$ Pressed $\to$ RT-Released) with rolling extrema. After actuation, release is emitted when depth falls by at least RT sensitivity below the running peak. From RT-Released, re-actuation occurs when depth rises by at least RT sensitivity above the running valley. The state resets to Idle only when the signal returns to zero, matching firmware reset semantics.

To preserve firmware behavior while keeping search efficient, signals are quantized before evaluation and searched on a discrete, hardware-relevant threshold grid.

\subsection{Search Space and Objective}
Optimization is discrete and exhaustive over the candidate threshold grid, yielding deterministic and reproducible solutions at exposed hardware resolution.

For each candidate, simulated click intervals are matched to labeled intervals using \textbf{overlap-constrained monotonic matching}. A label is matched only if at least one simulated interval overlaps that label interval. If multiple simulated intervals overlap the same label (oscillation/chatter), they are collapsed into one canonical matched interval using earliest make and latest break.

Let
\[
\begin{array}{rcl}
N  & = & \text{number of labeled clicks},\\
TP & = & \text{number of matched labels},\\
FN & = & N-TP,\\
FP & = & \text{number of unmatched simulated intervals}.
\end{array}
\]

For matched pair $i$, define make and break boundary errors:
\[
\begin{aligned}
e_i^{(m)} &= t_{i,\mathrm{sim}}^{(m)} - t_{i,\mathrm{label}}^{(m)},\\
e_i^{(b)} &= t_{i,\mathrm{sim}}^{(b)} - t_{i,\mathrm{label}}^{(b)}.
\end{aligned}
\]

Early errors are asymmetrically penalized with multiplier $\kappa \ge 1$:
\[
\phi_\kappa(e)=|e|\left(1+(\kappa-1)\mathbf{1}[e<0]\right).
\]

Per-match timing distance is
\[
d_i = w_{\mathrm{make}}\,\phi_\kappa\!\left(e_i^{(m)}\right)
+ \left(1-w_{\mathrm{make}}\right)\phi_\kappa\!\left(e_i^{(b)}\right),
\]
where $w_{\mathrm{make}} \in [0,1]$ controls make-vs-break priority.

Detection penalties are parameterized by base ms-equivalent detection cost $c_{\mathrm{det}}$ and recall/precision bias $\beta$:
\[
c_{\mathrm{miss}} = c_{\mathrm{det}}e^{\beta},\qquad
c_{\mathrm{fp}} = c_{\mathrm{det}}e^{-\beta}.
\]

The final loss is
\[
\mathcal{L}=
\begin{cases}
\infty, & TP=0,\\[4pt]
\displaystyle \frac{1}{TP}\sum_{i=1}^{TP} d_i
+ \frac{c_{\mathrm{miss}}FN + c_{\mathrm{fp}}FP}{N}, & TP>0.
\end{cases}
\]

This yields interpretable controls: $w_{\mathrm{make}}$ sets make-vs-break timing emphasis, $\kappa$ controls early-click severity, and $(c_{\mathrm{det}},\beta)$ control overall detection strictness and recall-precision preference.

\subsection{Pareto Frontier and Output}
For each feature-pair configuration, the optimizer returns all evaluated candidates, the minimum-loss operating point, and a Pareto frontier over \{error rate, average click distance\} for interactive inspection.

Detailed serialization choices (grid layout, sentinels, and per-combination diagnostics) are implementation concerns and are described in Section~\ref{sec:optimizer_impl}.
